\chapter{Mở Bài Bằng Cặp Đôi}
\chapterintro{Hai người nói trước, cả lớp nói sau}{Làm theo cặp là cách rất hiệu quả để kéo những lớp còn ngại phát biểu vào tiết học.}

\section{Quay sang bên cạnh nói một câu}
\begin{khoidongbox}{Hoạt động khởi động}
Cho học sinh quay sang bạn bên cạnh và nói đúng một câu về điều mình nhớ nhất hoặc đoán trước nhất.
\end{khoidongbox}
\begin{visaohieuqua}
Nói với một bạn dễ hơn nói trước cả lớp. Sau một lượt ngắn, lớp thường sẵn sàng phát biểu công khai hơn.
\end{visaohieuqua}
\begin{cachlam5phut}
Cho 20 giây trao đổi, rồi mời 2 cặp chia sẻ lại một ý. Không để hoạt động kéo dài thành thảo luận lớn.
\end{cachlam5phut}
\ngancach

\section{Hỏi bạn điều còn nhớ}
\begin{khoidongbox}{Hoạt động khởi động}
Mỗi em đặt cho bạn bên cạnh một câu rất ngắn: \textit{``Cậu còn nhớ gì từ buổi trước?''}
\end{khoidongbox}
\begin{visaohieuqua}
Khởi động bằng hỏi nhau làm lớp nhập vai người học ngay. Nó cũng biến bài cũ thành chuyện của cả hai, không chỉ của người bị gọi.
\end{visaohieuqua}
\begin{cachlam5phut}
Sau một lượt hỏi đáp, chọn 2 ý hay đưa lên bảng làm mốc vào bài mới.
\end{cachlam5phut}
\ngancach

\section{Cặp đôi đoán ý chính}
\begin{khoidongbox}{Hoạt động khởi động}
Trước khi học, cho từng cặp đoán xem ý chính của bài hôm nay sẽ là gì.
\end{khoidongbox}
\begin{visaohieuqua}
Đoán trước làm học sinh muốn kiểm chứng trong lúc học. Nhờ đó sự chú ý được kéo dài tốt hơn.
\end{visaohieuqua}
\begin{cachlam5phut}
Mỗi cặp chỉ nói một ý. Cuối tiết có thể quay lại xem cặp nào đoán gần nhất.
\end{cachlam5phut}
\ngancach

\section{Mỗi cặp tạo một ví dụ}
\begin{khoidongbox}{Hoạt động khởi động}
Cho hai người cùng nghĩ ra một ví dụ gần lớp mình nhất để hiểu bài hoặc chủ đề sắp học.
\end{khoidongbox}
\begin{visaohieuqua}
Tự tạo ví dụ là cách nhập bài rất mạnh. Làm theo cặp giúp áp lực nhẹ hơn và tốc độ nhanh hơn làm cá nhân.
\end{visaohieuqua}
\begin{cachlam5phut}
Gọi 3 cặp đọc ví dụ. Chọn một ví dụ gần nhất để dùng xuyên suốt trong tiết.
\end{cachlam5phut}
\ngancach

\section{Đổi bạn chốt cho nhau}
\begin{khoidongbox}{Hoạt động khởi động}
Mỗi em nói một điều ngắn, rồi bạn bên cạnh phải chốt lại bằng một câu rõ hơn.
\end{khoidongbox}
\begin{visaohieuqua}
Hoạt động này buộc học sinh nghe nhau thật sự chứ không chỉ nói cho xong. Nó rất hợp để tạo nhịp học chủ động đầu tiết.
\end{visaohieuqua}
\begin{cachlam5phut}
Chọn 2 cặp làm mẫu trước lớp, sau đó chốt: \textit{``Hôm nay mình sẽ học theo đúng tinh thần đó: nói và nghe cho rõ.''}
\end{cachlam5phut}